\documentclass[11pt,a4paper]{article}

\usepackage[T1]{fontenc}
\usepackage[utf8]{inputenc}
\usepackage[french]{babel}
\usepackage[margin=2.5cm]{geometry}
\usepackage{booktabs}
\usepackage{siunitx}
\usepackage[table]{xcolor}
\usepackage{hyperref}
\usepackage{pgfplots}
\pgfplotsset{compat=1.18}
\usepackage{caption}
\usepackage{enumitem}
\usepackage{amsmath}
\usepackage{amssymb}

\sisetup{detect-weight=true, detect-family=true}

\hypersetup{
  colorlinks=true,
  linkcolor=blue!50!black,
  urlcolor=blue!50!black,
  citecolor=blue!50!black,
  pdftitle={Note - Enthalpie relative Al4C3},
  pdfauthor={Pr. John Akapulco}
}

\title{\vspace{-1.5cm}Note technique\\[2pt]
\large Diagramme d'enthalpie relative $\Delta H(P)$ d'Al$_4$C$_3$ (référence R-3m)}
\author{Pr. John Akapulco \\ \small \href{mailto:computchem86@gmail.com}{computchem86@gmail.com}}
\date{23 août 2026}

\begin{document}
\maketitle
\vspace{-0.5cm}

\begin{abstract}
\noindent
Cette note complète le rapport \texttt{rapport\_activite\_Al4C3\_HP\_scan.tex} (scan H(P)
complet, 0--100~GPa, sept polymorphes) en présentant les mêmes résultats sous forme
d'enthalpie \emph{relative} : $\Delta H(P) = H_{\text{phase}}(P) - H_{\text{R-3m}}(P)$, exprimée
en meV/atome, où \textbf{R-3m} est prise comme référence et reste fixée à $\Delta H = 0$ sur
toute la gamme de pression. Cette normalisation rend les écarts de stabilité directement
lisibles en meV/atome. Avec R-3m comme référence, \textbf{Pnma-VII} reste en dessous de zéro
sur \emph{toute} la gamme étudiée (de $-3.441$~meV/atome à 0~GPa à $-28.689$~meV/atome à
100~GPa) : R-3m n'est jamais la phase la plus stable des sept candidats, mais elle n'en est
séparée que de quelques meV/atome à basse pression. Les autres phases (Pbam, Pnma-III, Pnma,
anti-CaFe$_2$O$_4$, I-43d) démarrent nettement au-dessus de zéro à 0~GPa et croisent la
référence successivement entre $\approx$23 et $\approx$45~GPa à mesure que la pression
augmente.
\end{abstract}

\section{Construction du diagramme}

Pour chaque polymorphe $i$ et chaque pression $P$ de la grille (0 à 100~GPa, pas de 10~GPa),
on calcule
\[
\Delta H_i(P) = H_i(P) - H_{\text{R-3m}}(P), \qquad
\Delta H_i(P)\ \text{en meV/atome} = \frac{\Delta H_i(P)\ [\text{eV/f.u.}]}{7}\times 1000,
\]
avec 7 atomes par formule unitaire (Al$_4$C$_3$ = 4~Al + 3~C). R-3m est retenue ici comme phase
de référence ; ce n'est pas la phase de plus basse enthalpie à 0~GPa dans les sept candidats
(voir tableau~\ref{tab:rank0}) -- Pnma-VII la devance de 0.024~eV/f.u.
($\approx$3.441~meV/atome), un écart proche de l'incertitude numérique du maillage
$\mathbf{k}$-points allégé utilisé pour ce scan (cf.\ rapport principal, section
« Robustesse du classement »). Soustraire la même courbe $H_{\text{ref}}(P)$ à chaque phase ne
déplace aucune des pressions de transition rapportées dans le rapport principal -- ces
transitions dépendent uniquement de la différence entre deux courbes, invariante par un
changement de référence commun ; seule change la position de la droite $\Delta H = 0$.

\begin{table}[h]
\centering
\begin{tabular}{clS[table-format=-2.3]}
\toprule
Rang & Polymorphe & {$H$/f.u. à 0~GPa (eV)} \\
\midrule
1 & Pnma-VII           & -43.354 \\
2 & R-3m (référence)   & -43.329 \\
3 & Pnma-III           & -42.360 \\
4 & Pbam               & -42.226 \\
5 & Pnma               & -41.596 \\
6 & anti-CaFe$_2$O$_4$ & -41.502 \\
7 & I-43d              & -40.232 \\
\bottomrule
\end{tabular}
\caption{Classement à 0~GPa (jeu complet, 7/7 convergés), valeurs au millième d'eV. R-3m, prise
comme référence pour le diagramme relatif, est la deuxième phase la plus stable à 0~GPa, à
0.024~eV/f.u.\ de Pnma-VII.}
\label{tab:rank0}
\end{table}

\section{Diagramme d'enthalpie relative}

\begin{figure}[h]
\centering
\begin{tikzpicture}
\begin{axis}[
  width=\textwidth, height=8.5cm,
  xlabel={Pression (GPa)},
  ylabel={$\Delta H$/atome (meV), réf. R-3m},
  enlarge y limits={upper=0.22, lower=0.02},
  legend style={
    at={(0.5,0.98)}, anchor=north,
    legend columns=4,
    font=\scriptsize,
    fill=white, fill opacity=0.85, draw opacity=1, text opacity=1,
    draw=black!20,
    /tikz/every even column/.append style={column sep=6pt},
  },
  grid=major,
  extra y ticks={0},
  extra y tick style={grid style={solid, black!40}},
]
\addplot table[x=P_GPa,y=dH_meV_per_atom] {../../hp_curves/hp_rel_R-3m.dat};
\addplot table[x=P_GPa,y=dH_meV_per_atom] {../../hp_curves/hp_rel_Pnma-VII.dat};
\addplot table[x=P_GPa,y=dH_meV_per_atom] {../../hp_curves/hp_rel_Pbam.dat};
\addplot table[x=P_GPa,y=dH_meV_per_atom] {../../hp_curves/hp_rel_Pnma-III.dat};
\addplot table[x=P_GPa,y=dH_meV_per_atom] {../../hp_curves/hp_rel_Pnma.dat};
\addplot table[x=P_GPa,y=dH_meV_per_atom] {../../hp_curves/hp_rel_anti-CaFe2O4.dat};
\addplot table[x=P_GPa,y=dH_meV_per_atom] {../../hp_curves/hp_rel_I-43d.dat};
\legend{R-3m (réf.), Pnma-VII, Pbam, Pnma-III, Pnma, anti-CaFe2O4, I-43d}
\end{axis}
\end{tikzpicture}
\caption{Enthalpie relative $\Delta H(P) = H_{\text{phase}}(P) - H_{\text{R-3m}}(P)$, en
meV/atome. La phase de référence est la droite $\Delta H = 0$ ; toute courbe qui passe sous zéro
devient, à cette pression, plus stable que R-3m. Pnma-VII reste sous zéro sur toute la gamme.
Les cinq autres phases démarrent au-dessus de zéro à 0~GPa et croisent la référence entre
$\approx$23 et $\approx$45~GPa (voir section~\ref{sec:lecture}).}
\label{fig:hprel}
\end{figure}

\section{Lecture du diagramme}
\label{sec:lecture}

\begin{itemize}
  \item \textbf{Pnma-VII} est la seule courbe à rester négative sur toute la gamme : de
  $-3.441$~meV/atome à 0~GPa à $-28.689$~meV/atome à 100~GPa. R-3m n'est donc jamais la phase la
  plus stable des sept, mais l'écart à 0~GPa (3.4~meV/atome) est faible.
  \item \textbf{0~GPa} : hors Pnma-VII, toutes les autres phases sont nettement au-dessus de
  zéro -- Pnma-III (+138.504), Pbam (+157.682), Pnma (+247.622), anti-CaFe$_2$O$_4$ (+261.015)
  et I-43d (+442.527)~meV/atome.
  \item \textbf{Croisements avec la référence} (interpolation linéaire entre les deux points de
  grille encadrant le changement de signe) :
  \begin{center}
  \begin{tabular}{lc}
  \toprule
  Polymorphe & Pression de croisement (vs R-3m) \\
  \midrule
  Pnma-III           & 23.4~GPa \\
  Pbam               & 23.9~GPa \\
  Pnma               & 30.2~GPa \\
  anti-CaFe$_2$O$_4$ & 34.3~GPa \\
  I-43d              & 45.3~GPa \\
  \bottomrule
  \end{tabular}
  \end{center}
  Pbam et Pnma-III dépassent R-3m en stabilité quasi simultanément, un peu avant la transition
  Pnma-VII~$\to$~Pbam (25.6~GPa) identifiée dans le rapport principal sur l'enveloppe de
  stabilité absolue.
  \item \textbf{I-43d} reste la courbe la plus haute jusqu'à $\approx$45~GPa (elle passe sous
  zéro entre 40 et 50~GPa), avec l'écart à la référence le plus marqué à 0~GPa
  (+442.527~meV/atome), signe d'une convergence progressive des enthalpies sous forte
  compression.
\end{itemize}

\appendix
\section{Références}

Données sources : \texttt{hp\_curves/results\_hp.csv} (commit \texttt{470e557}, campagne HP
scan 0--100~GPa). Fichiers de données relatifs générés pour cette note :
\texttt{hp\_curves/hp\_rel\_<polymorphe>.dat} (colonnes \texttt{P\_GPa},
\texttt{dH\_meV\_per\_atom}, référence R-3m). Voir \texttt{rapport\_activite\_Al4C3\_HP\_scan.tex}
pour la méthodologie DFT complète, les diagrammes $V(P)$/$E(P)$/$H(P)$ en valeurs absolues et le
tableau numérique complet (E, V, H).

\end{document}
